% Appendix O (v3): prompt and defect architecture. The v2 appendix printed every logged exchange in
% full (about 30 pages); v3 keeps one table of every exhibited stage, the two exchanges a reader
% needs to reproduce the method, and the documented defects as callouts. The complete exchanges
% are unchanged in docs/paper/appendix/k-prompts.tex (v1/v2) and in the released call logs;
% excerpts below are copied from them verbatim (scripts/papers/prompt_appendix.py).
\clearpage
\onecolumn
\begingroup
\newtcblisting{nwprompt}[3][]{
  enhanced, breakable, listing only,
  colback=#2!8!white, colframe=#2,
  colbacktitle=#2!35!white, coltitle=nwink,
  title={#3}, fonttitle=\sffamily\bfseries\footnotesize,
  boxrule=0.45pt, arc=1.5mm, left=1.5mm, right=1.5mm, top=1mm, bottom=1mm,
  before skip=5pt, after skip=6pt,
  listing options={basicstyle=\ttfamily\scriptsize, columns=fullflexible, keepspaces=true,
    breaklines=true, breakatwhitespace=false, breakindent=1em, showstringspaces=false,
    frame=none, aboveskip=0pt, belowskip=0pt},
  #1}
\newtcolorbox{nwdefect}[1]{enhanced, breakable, colback=nwhighpink!45!white, colframe=nwMagenta,
  boxrule=0.5pt, arc=1.5mm, left=2mm, right=2mm, top=1mm, bottom=1mm, before skip=6pt,
  after skip=6pt, title={#1}, fonttitle=\sffamily\bfseries\footnotesize, colbacktitle=nwmagenta,
  coltitle=nwink, fontupper=\small}

\section{Prompts and Model Exchanges}
\label{app:prompts}
Every generative stage calls the generator through one client that records the system prompt,
the user prompt and the raw response. Table~\ref{tab:prompt-stages} lists the stages whose
first successful call from the run behind the reported results is exhibited in the release;
all used \res{GEN-model}. It first prints the two exchanges a reader needs to reproduce the method (same-volume
arbitration and the paraphrase that makes the treated text) and the five defects that the
logged prompts and responses carry; the complete logged exchange of every stage follows
(Appendix~\ref{app:prompts-full}). Every call log is released
(Appendix~\ref{app:reproducibility}); the excerpts are verbatim, with elisions marked. Evaluation labels were not produced by
these stages (Appendix~\ref{app:annotation}).

\begin{table*}[t]
\centering\footnotesize
\setlength{\tabcolsep}{4pt}
\begin{tabular}{@{}>{\raggedright\arraybackslash}p{3.0cm}>{\raggedright\arraybackslash}p{6.3cm}r>{\raggedright\arraybackslash}p{3.75cm}@{}}
\toprule
\nwhead{Stage} & \nwhead{Role} & \nwhead{Tokens in / out} & \nwhead{Run} \\
\midrule
\multicolumn{4}{@{}l}{\emph{\sys build}} \\
Entity typing & Types each mined name candidate (gazetteer) (\emph{``Anne'', v1}) & 518 / 39 & \nolinkurl{20261004T193719-gazetteer-v1-5} \\
Alias adjudication & Decides whether two surface forms are one entity (gazetteer) (\emph{Charlie vs.\ Carrie Sloane}) & 548 / 50 & \nolinkurl{20261004T193719-gazetteer-v1-5} \\
Claim extraction & Extracts dated, quoted claims about one character from one evidence window (\emph{Anne Shirley, v5}) & 8,884 / 734 & \nolinkurl{20261005T070204-extract-v5-5} \\
Scene extraction & Records scenes, participants and events over a chapter span (\emph{v5, chapter 2}) & 5,729 / 202 & \nolinkurl{20261004T195945-scenes-v5-5} \\
Infobox pass & Long-context pass for gender, age and origin (\emph{Anne Shirley, v1}) & 93,403 / 490 & \nolinkurl{20261005T070216-infobox-v1-5} \\
Backstory pass & Long-context pass for background facts (\emph{Marilla Cuthbert, v1}) & 135,267 / 309 & \nolinkurl{20261005T070229-backstory-v1-5} \\
Arbitration & Resolves a same-volume conflict between two values (\emph{Anne Shirley, AGE, v1}) & 258 / 55 & \nolinkurl{20261005T070243-graph_build-v1-5} \\
Verification & Checks a stored fact against its own cited evidence (\emph{Anne Shirley, APPEARANCE}) & 2,833 / 5 & \nolinkurl{20261005T053635-verify-upto5} \\
\midrule
\multicolumn{4}{@{}l}{\emph{\sys publication}} \\
Page prose (appearance) & Writes a physical description from the cutoff-filtered claims (\emph{Anne Shirley, $t{=}1$}) & 1,045 / 174 & \nolinkurl{20261005T070347-synthesize-upto5} \\
Codex summary & Writes the one-sentence summary of a shared codex entry (\emph{Aid Society, $t{=}1$}) & 1,974 / 69 & \nolinkurl{20261005T070527-site_build-upto5} \\
\midrule
\multicolumn{4}{@{}l}{\emph{Comparison systems and diagnostics}} \\
B1 page & Long context: one character page from the whole prefix (\emph{Anne Shirley, $t{=}4$}) & 559,824 / 5,688 & \nolinkurl{20261001T043209-baseline-b1-v4} \\
B2 page & Prefix retrieval: one character page from retrieved passages (\emph{Marilla Cuthbert, $t{=}1$}) & 24,083 / 2,477 & \nolinkurl{20261002T095544-baseline-b2-v1} \\
B3 index & LightRAG's own entity and relation extraction while building its index & 3,243 / 1,022 & \nolinkurl{20260929T225346-baseline-b3-v5} \\
B3 page & LightRAG: one character page from LightRAG context (\emph{Anne Shirley, $t{=}5$}) & 17,438 / 2,191 & \nolinkurl{20261001T142328-baseline-b3-v5} \\
X1 page & Closed book: title and character name only (\emph{Matthew Cuthbert, $t{=}5$}) & 451 / 1,087 & \nolinkurl{20260930T214444-baseline-x1-v5} \\
X2 page & Full text: all five volumes and the cutoff instruction (\emph{Gilbert Blythe, $t{=}1$}) & 689,967 / 2,449 & \nolinkurl{20260930T222548-baseline-x2-v1} \\
\midrule
\multicolumn{4}{@{}l}{\emph{Treated text and recognition}} \\
Paraphrase & Paraphrases one renamed paragraph & 385 / 209 & \nolinkurl{20261003T165238-decontaminate-v5} \\
Memorization probe & Name cloze or title identification on one paragraph & 227 / 9 & \nolinkurl{20261003T165318-probe-memorization-v1-5} \\
\bottomrule
\end{tabular}
\caption{Every exhibited model stage: the first successful call of the run behind the reported
results, with the exhibited character or entry and volume or cutoff in italics. Runs are under \nolinkurl{anne@v2/} (build and pages) or \nolinkurl{anne-decon/}
(treatment and probe) in the release. B4's page prompt, the other prose sections and the
shared-page summary check are in the release only.}
\label{tab:prompt-stages}
\end{table*}

\subsection*{Same-volume arbitration}
Two values of one attribute from one volume are treated as an extraction conflict
(Section~\ref{sec:method}). A \texttt{flag} resolution keeps the higher-confidence value
provisionally (here 11), so the interval sequence has no gap, and records the conflict.
\begin{nwprompt}{nwViolet}{User prompt (system prompt in the release)}
Character: anne-shirley
Attribute: AGE
Volume: 1

Fact A: value="11", confidence=1.00, evidence: "I was eleven last March," said Anne"
Fact B: value="13", confidence=0.99, evidence: "It's such a solemn thing to be almost fourteen, Marilla."; ""Just think, Diana, I'm thirteen years old today," remarked Anne in an awed voice."

Which fact is correct? resolution must be one of keep_a, keep_b, keep_both, flag, plus a one-sentence rationale citing the evidence.
\end{nwprompt}
\begin{nwprompt}{nwCyan}{Model output}
{"resolution": "flag",
 "rationale": "Fact A states Anne was eleven based on an earlier passage, while Fact B shows her turning thirteen later in the same volume, indicating a progression of time rather than a single static age."}
\end{nwprompt}

\subsection*{Paraphrase for the treated text}
Each renamed paragraph is paraphrased in one call, keeping its identifier.
\begin{nwprompt}{nwViolet}{System and user prompt (paragraph abridged)}
SYSTEM: You rewrite passages of fiction in new words for a research corpus. You never add, remove or explain anything.

Rewrite the paragraph below in different words and different sentence structure.
Rules:
- Keep every fact, event, number, date and place, and keep them in the same order.
- Keep every personal and ship name exactly as written. Do not introduce any name that is not in the paragraph.
- Dialogue stays dialogue, spoken by the same person. Use the same quotation marks as the original (" " stay " "). Only spoken words go inside them; a speech tag such as "said I" stays outside.
- Keep the narrator's person (first person stays first person).
- Do not summarise, shorten substantially, or add commentary.
Return only the rewritten paragraph, nothing else.

Paragraph:
It was a clear, apple-green evening in May, and Seven Gulls Harbour was mirroring back the clouds of the golden west between its softly dark shores. [...]
\end{nwprompt}
\begin{nwprompt}{nwCyan}{Model output (abridged)}
The May evening was transparent and apple-green, as Seven Gulls Harbour reflected the clouds from the golden west amidst its gently darkened banks. [...]
\end{nwprompt}

\subsection*{Defects carried by the logged prompts and responses}
The five defects below are part of the reported runs. They are documented, not repaired;
Section~\ref{sec:discussion} and Appendix~\ref{app:formal} relate each to the results.

\begin{nwdefect}{D1\enspace Implementation notes used as field definitions (extraction and verification)}
The extraction prompt defines STATUS and APPEARANCE with internal revision labels and code
paths, and the verifier is told to judge an APPEARANCE fact by that same text:
\begin{nwprompt}{nwMagenta}{Verification user prompt (excerpt)}
Fact: APPEARANCE = "little pale face with seven freckles on the nose" -- evidence: [v02:c27:p0041] [...]
Definition of APPEARANCE: subsume (Phase 22 A2): drop a value that is redundant next to a longer sibling ('hair' next to 'long brown hair') or a bare single-word noun standing alone -- see graph/canonicalize.py::subsume_values. Deterministic, no embed call.
\end{nwprompt}
The exhibited claim is accepted; the example alone does not measure the defect's effect.
\end{nwdefect}

\begin{nwdefect}{D2\enspace Responses depart from prompt rules the code does not enforce (infobox pass)}
The volume-1 infobox response gives a \texttt{GENDER} quote that does not contain Anne's
name, and two \texttt{AGE} values (11 and 13). The code applies the evidence gates of
Section~\ref{sec:method}, not these two prompt rules; the two ages go to arbitration, shown
above.
\end{nwdefect}

\begin{nwdefect}{D3\enspace Early-volume prompts carry later aliases}
Alias lists come from the whole-series entity index, so the volume-1 infobox call already
names Anne by her married names:
\begin{nwprompt}{nwMagenta}{Infobox user prompt, volume 1 (excerpt)}
Character: "Anne Shirley" (aliases: Anne, Anne-girl, Mistress Blythe, Mrs. Blythe, Mrs. Doctor, Mrs. Doctor, dear, Mrs. Dr. dear, Queen Anne, the Anne-girl, young Mrs. Doctor)
Volume 1.
\end{nwprompt}
This is construction dependence reaching the generator's input
(Section~\ref{sec:results-scope}); the comparisons of this paper do not isolate it.
\end{nwdefect}

\begin{nwdefect}{D5\enspace A response cites a paragraph ID where a line key was requested (codex summary)}
The codex prompt asks the model to cite input-line keys such as \texttt{W1}; the response
returns \texttt{v01:c27:p0000}, a paragraph ID. The parser normalizes returned paragraph IDs
during citation resolution, which maps an identifier to a paragraph but does not check that
the paragraph supports the sentence.
\end{nwdefect}

\begin{nwdefect}{D4\enspace Residue of one prompt serving every corpus}
The typing prompt lists a type unused here (\texttt{TECH}) and, as logged, repeats one example
sentence three times; the extraction prompt's JSON example uses a character from another
series (``Shin''); the paraphrase rule keeps ``personal and ship'' names, which has no effect
on this series.
\end{nwdefect}
\endgroup
